%% LaRT Clump Medium Technical Description
%% Author: Kwang-Il Seon
\documentclass[english,a4paper]{kiseon_memo}
\synctex=1
\usepackage{listings}
\usepackage{xcolor}
\usepackage{amsmath}
\usepackage{amssymb}

\lstset{
  basicstyle=\ttfamily\small,
  keywordstyle=\color{blue}\bfseries,
  commentstyle=\color{gray},
  stringstyle=\color{olive},
  frame=single,
  breaklines=true,
  numbers=left,
  numberstyle=\tiny\color{gray},
  language=[95]Fortran
}

\usepackage{babel}
\begin{document}

\title{Clumpy Medium Extension of LaRT: Algorithm Description}
\author{Kwang-Il Seon}
\date{2026-04-24}
\maketitle

\tableofcontents
\bigskip

%=============================================================================
\section{Overview}
%=============================================================================

The clumpy medium mode adds a population of $N$ non-overlapping spherical
clumps inside a spherical domain of radius $R$.  Each clump is a sphere of
radius $r_c$ with uniform Lyman-$\alpha$ opacity; the inter-clump medium is
vacuum.

The mode is activated by
\begin{lstlisting}
par%use_clump_medium = .true.
\end{lstlisting}
in the input namelist.  All raytrace procedure pointers are redirected to
clump-aware routines (\texttt{raytrace\_clump.f90}) while scatter and
peel-off reuse the standard Cartesian routines from
\texttt{scattering\_car.f90} and \texttt{peelingoff\_rect.f90}, since the
uniform clump physical parameters ($T$, $v_\mathrm{bulk}$) are broadcast to
the Cartesian grid arrays ($\texttt{grid\%Dfreq}$, $\texttt{grid\%voigt\_a}$).

New source files introduced for this mode:
\begin{center}
\begin{tabular}{ll}
\hline
File & Role \\
\hline
\texttt{clump\_mod.f90}       & Clump population: init, RSA placement, CSR grid, raytrace helpers \\
\texttt{raytrace\_clump.f90}  & Raytrace through the clump medium (DDA + sphere intersection) \\
\texttt{grid\_mod\_clump.f90} & Grid setup: creates Cartesian shell grid, then calls \texttt{init\_clumps} \\
\hline
\end{tabular}
\end{center}

%=============================================================================
\section{Physical Setup}
%=============================================================================

\subsection{Geometry}

The outer domain is a sphere of radius $R = \texttt{par\%rmax}$ (code units).
$N$ spherical clumps of radius $r_c = \texttt{par\%clump\_radius}$ are placed
at random inside this sphere (non-overlapping).

A Cartesian grid of $N_x \times N_y \times N_z$ cells is created with the
same bounding box $[-R, R]^3$.  Grid arrays ($\texttt{rhokap}$, $\texttt{Dfreq}$,
$\texttt{voigt\_a}$, $\texttt{vf}$) are set to uniform values:
\begin{itemize}
  \item $\texttt{rhokap} = 0$ everywhere (opacity lives only inside clumps);
  \item $\texttt{Dfreq} = \Delta\nu_D^{cl}$ (clump Doppler frequency);
  \item $\texttt{voigt\_a} = a_V^{cl}$ (clump Voigt parameter);
  \item $\texttt{vfx} = \texttt{vfy} = \texttt{vfz} = 0$.
\end{itemize}
Peel-off and scattering routines read these uniform grid values, so no
special-casing is needed inside those modules.

\subsection{Clump Count}

The user specifies one of three equivalent parameters; the code derives $N$
from whichever is set:

\paragraph{From $N_\mathrm{cov}$ (covering factor):}
The mean number of clumps intercepted by a radial ray from the centre to the
sphere surface is
\begin{equation}
  N_\mathrm{cov} = N\,\frac{\pi r_c^2}{\pi R^2}\cdot\frac{R}{\frac{4}{3}\pi R^3/(4\pi R^2)}
  = \frac{3}{4}\,N\!\left(\frac{r_c}{R}\right)^2 ,
\end{equation}
hence
\begin{equation}
  N = \frac{4}{3}\,N_\mathrm{cov}\!\left(\frac{R}{r_c}\right)^2 .
\end{equation}
Input parameter: \texttt{par\%clump\_N\_cov}.

\paragraph{From $f_\mathrm{vol}$ (volume filling factor):}
\begin{equation}
  f_\mathrm{vol} = N\left(\frac{r_c}{R}\right)^3 ,
  \qquad N = f_\mathrm{vol}\!\left(\frac{R}{r_c}\right)^3 .
\end{equation}
Input parameter: \texttt{par\%clump\_f\_vol}.

\paragraph{Direct specification:} \texttt{par\%clump\_N\_clumps}.

\subsection{Clump Physical Parameters}

The clump temperature $T_c$ defaults to $\texttt{par\%temperature}$ unless
\texttt{par\%clump\_temperature} $> 0$.  Derived quantities:
\begin{align}
  v_\mathrm{th} &= v_{\mathrm{th},1}\sqrt{T_c},\quad
                   v_{\mathrm{th},1} = \sqrt{2k_B/m_H}\approx 12.84\;\mathrm{km\,s^{-1}/\sqrt{10^4\,K}},\\
  \Delta\nu_D^{cl} &= v_\mathrm{th} / (\lambda_0 \times \mathrm{km2um}),\\
  a_V^{cl} &= \frac{A_{21}}{4\pi\,\Delta\nu_D^{cl}},
\end{align}
where $\lambda_0$ is the Lyman-$\alpha$ rest wavelength.

The line-centre optical depth of a clump (centre to surface) is
\texttt{par\%clump\_tau0}; from this the uniform opacity is
\begin{equation}
  \kappa_c = \frac{\tau_0}{\phi(0,a_V^{cl})\, r_c},
\end{equation}
where $\phi(x, a)$ is the Voigt profile at dimensionless frequency $x$.
Alternatively, the hydrogen number density \texttt{par\%clump\_nH} may be
given, in which case $\kappa_c = n_H\,\sigma_0\,d_\mathrm{cm}$ (requires
\texttt{par\%distance\_unit}).

\subsection{Clump Velocities}

Each clump's velocity is the sum of two independent components:
\begin{equation}
  \mathbf{v}_i = \mathbf{v}_\mathrm{sys}(\mathbf{r}_i) + \mathbf{v}_\mathrm{rand,i},
\end{equation}
where $\mathbf{v}_\mathrm{sys}$ is a position-dependent systematic flow
(set by \texttt{par\%velocity\_type}) and $\mathbf{v}_\mathrm{rand}$ is an
independent Gaussian random component (set by \texttt{par\%clump\_sigma\_v}).
Either component can be zero.

\paragraph{Random component:}
If \texttt{par\%clump\_sigma\_v} $> 0$, each clump is assigned
$\mathbf{v}_{\mathrm{rand},i} \sim \mathcal{N}(0,\sigma_v^2)$ per component
(drawn in \texttt{generate\_clumps}).

\paragraph{Systematic component (\texttt{par\%velocity\_type}):}
After RSA placement, \texttt{assign\_clump\_velocities\_from\_type} evaluates
the flow velocity at each clump centre and \emph{adds} it to the existing
(possibly non-zero) velocity.  Supported types:

\begin{center}
\begin{tabular}{lll}
\hline
\texttt{velocity\_type} & Formula & Parameters \\
\hline
\texttt{'hubble'} & $\mathbf{v} = V_\mathrm{exp}\,\mathbf{r}/R$ & \texttt{Vexp} \\
\texttt{'constant\_radial'} & $\mathbf{v} = V_\mathrm{exp}\,\hat{\mathbf{r}}$ & \texttt{Vexp} \\
\texttt{'parallel\_velocity'} & $\mathbf{v} = (V_x, V_y, V_z)$ & \texttt{Vx, Vy, Vz} \\
\texttt{'ssh'} & see below & \texttt{Vpeak, rpeak, DeltaV} \\
\texttt{'rotating\_solid\_body'} & $\mathbf{v} = V_\mathrm{rot}(-y,x,0)/R$ & \texttt{Vrot} \\
\texttt{'rotating\_galaxy\_halo'} & flat rotation about $z$ & \texttt{Vrot, rinner} \\
\hline
\end{tabular}
\end{center}

All velocity magnitudes are in km/s.  The \texttt{'ssh'} model (Song, Seon \&
Hwang 2020) gives a linear rise inside $r_\mathrm{peak}$ and a Hubble-like
outer flow:
\begin{equation}
  v(r) = \begin{cases}
    V_\mathrm{peak}\,r / r_\mathrm{peak} & r < r_\mathrm{peak}, \\
    V_\mathrm{peak} + \Delta V\,(r - r_\mathrm{peak})/(R - r_\mathrm{peak}) & r \ge r_\mathrm{peak}.
  \end{cases}
\end{equation}
The \texttt{'rotating\_galaxy\_halo'} model uses the cylindrical radius
$\rho = \sqrt{x^2+y^2}$ and gives solid-body rotation inside $r_\mathrm{inner}$
and a flat curve ($v_\phi = V_\mathrm{rot}$) outside.

%=============================================================================
\section{Clump Placement: Random Sequential Addition (RSA)}
%=============================================================================

\subsection{Algorithm}

Clumps are placed one at a time.  Each candidate position is drawn uniformly
inside the sphere (box-rejection sampling) and accepted only if it does not
overlap any previously placed clump (centre separation $\ge 2r_c$).

Overlap checking is accelerated by a linked-list hash grid of cell size
$\ell_\mathrm{RSA} \approx 2R/N_g$ where $N_g = \min(512, \lfloor N^{1/3}
\rfloor + 1)$.  Only the $3^3 = 27$ neighbouring cells of the candidate are
checked, reducing the expected check count to $O(1)$ per attempt (for
$f_\mathrm{vol} < 35\%$).

\begin{lstlisting}
do while (icl < N_clumps)
   !--- draw uniform point inside sphere
   do
      xc = (2*rand - 1)*R; yc = ...; zc = ...
      if (xc^2+yc^2+zc^2 <= R^2) exit
   enddo
   !--- check 27 neighbours for overlap (min_sep = 2*r_c)
   overlap = .false.
   outer: do kg2=...; do jg2=...; do ig2=...
      jnb = head(cell)
      do while (jnb > 0)
         if (dist2(xc,yc,zc, cl_x(jnb),...) < (2*r_c)^2) then
            overlap=.true.; exit outer
         endif
         jnb = nxt(jnb)
      enddo
   enddo; enddo; enddo outer
   if (overlap) cycle
   icl = icl+1
   cl_x(icl)=xc; cl_y(icl)=yc; cl_z(icl)=zc
   ! assign velocity if sigma_v > 0
   ! insert into linked list
enddo
\end{lstlisting}

\subsection{Acceptance Rate}

For the test case ($N_\mathrm{cov}=5$, $r_c/R=0.01$, $N=66667$,
$f_\mathrm{vol}=0.067$) the RSA acceptance rate is approximately 74\%.

%=============================================================================
\section{CSR Acceleration Grid}
%=============================================================================

After RSA placement, a compressed-sparse-row (CSR) grid is built to
accelerate ray--clump intersection queries during radiative transfer.

\subsection{Structure}

A 3-D Cartesian grid of $M_x \times M_y \times M_z$ cells ($M = N_g$ from above)
covers the domain $[-R-r_c, R+r_c]^3$.  Each clump is registered in every CSR
cell whose bounding box overlaps the clump sphere.  The CSR layout uses two
arrays:
\begin{itemize}
  \item \texttt{cg\_start(icell)} -- index into \texttt{cg\_list} of the first
        entry for cell \texttt{icell};
  \item \texttt{cg\_list(j)} -- clump index stored at position $j$.
\end{itemize}

\subsection{Construction (Two-Pass)}

Pass 1: count registrations per cell $\to$ prefix-sum to fill
\texttt{cg\_start}.\\
Pass 2: fill \texttt{cg\_list}.

A clump is registered in cell $(i,j,k)$ if the cell's AABB intersects the
clump sphere, i.e.\ the cell range in each dimension is
$[\lfloor(c_x-r_c-x_\mathrm{min})/\delta\rfloor,
  \lfloor(c_x+r_c-x_\mathrm{min})/\delta\rfloor]$.

Both \texttt{cg\_start} and \texttt{cg\_list} are allocated as MPI-3
shared-memory windows (one copy per node).

%=============================================================================
\section{Raytrace Through the Clump Medium}
%=============================================================================

\subsection{Overview}

Photon propagation is split into two phases:

\begin{enumerate}
  \item \textbf{Advance through vacuum} -- use the DDA traversal of the CSR
        grid (\S\ref{sec:dda}) to find the next clump hit.
  \item \textbf{Traverse through a clump} -- accumulate optical depth along
        a straight path through the sphere (\S\ref{sec:clump_traverse}).
\end{enumerate}

Frequency shifts due to individual clump bulk velocities are applied at
clump entry and reversed at clump exit.

\subsection{Ray--Sphere Intersection}

Given a ray $\mathbf{p}(t) = \mathbf{o} + t\hat{\mathbf{k}}$ and a sphere
centred at $\mathbf{c}$ with squared radius $r_c^2$, define
$\mathbf{r} = \mathbf{o} - \mathbf{c}$, $b = \mathbf{r}\cdot\hat{\mathbf{k}}$,
$\Delta = b^2 - |\mathbf{r}|^2 + r_c^2$.  If $\Delta < 0$ there is no
intersection.  Otherwise
\begin{equation}
  t_\mathrm{entry} = -b - \sqrt{\Delta},\qquad
  t_\mathrm{exit}  = -b + \sqrt{\Delta}.
\end{equation}
A hit is recorded when $t_\mathrm{exit} > 0$.

\subsection{DDA Search for Next Clump}
\label{sec:dda}

The Amanatides \& Woo (1987) digital differential analyser (DDA) traverses
the CSR cells along the ray in order of increasing path length.  For each
cell, all registered clumps are tested with the ray--sphere formula.  The
traversal stops as soon as the current cell's entry distance exceeds the
best hit distance found so far.

\begin{lstlisting}
best_te  = huge;  best_icl = 0
d        = 0     ! path length to current cell entry

do while (.true.)
   if (d > best_te .or. d > t_max) exit
   ! test all clumps in current cell
   do ip = cg_start(icell), cg_start(icell+1)-1
      icl = cg_list(ip)
      if (icl == skip_icl) cycle
      call ray_sphere_isect(..., te, tx2, hit)
      if (hit .and. tx2>0 .and. te<best_te) then
         best_te=te; best_tx2=tx2; best_icl=icl
      endif
   enddo
   ! advance to next cell (Amanatides & Woo)
   if (tx<=ty .and. tx<=tz) then
      d=tx; ci=ci+si; tx=tx+delx
   elseif (ty<=tz) then
      d=ty; cj=cj+sj; ty=ty+dely
   else
      d=tz; ck=ck+sk; tz=tz+delz
   endif
enddo
\end{lstlisting}

\subsection{Optical Depth Traversal}
\label{sec:clump_traverse}

\texttt{raytrace\_to\_tau\_clump} advances the photon until optical depth
$\tau_\mathrm{in}$ is accumulated or the outer sphere is exited.

\paragraph{Inside a clump ($\texttt{icell\_clump} > 0$):}
The remaining path through the current clump is
$s = t_\mathrm{exit}(\mathbf{p}, \hat{\mathbf{k}}, \mathrm{icl})$.
Opacity at the current frequency:
$\kappa = \kappa_c\,\phi(x, a_V)$.
If $\tau_\mathrm{rem} \le \kappa s$, scatter at $\mathrm{d}s = \tau_\mathrm{rem}/\kappa$.
Otherwise subtract $\kappa s$ from $\tau_\mathrm{rem}$, advance to clump exit,
shift frequency back to the global frame, and continue in vacuum.

\paragraph{In vacuum ($\texttt{icell\_clump} = 0$):}
Call \texttt{find\_next\_clump} (DDA, \S\ref{sec:dda}) to find the nearest
clump.  If none is found along the ray, advance to the sphere boundary and
exit ($\texttt{photon\%inside} = \texttt{.false.}$).

\paragraph{Frequency convention:}
The photon frequency $x$ (in units of $\Delta\nu_D^{cl}$) is maintained in
the local clump rest frame while inside a clump, and in the global frame in
vacuum:
\begin{align}
  \text{at clump entry:}\quad &x \leftarrow x - u_\mathrm{los}/v_\mathrm{th},\\
  \text{at clump exit:}\quad  &x \leftarrow x + u_\mathrm{los}/v_\mathrm{th},
\end{align}
where $u_\mathrm{los} = (\mathbf{v}_i \cdot \hat{\mathbf{k}})$.

\paragraph{Jout accumulation:}
When the photon exits the outer sphere (\texttt{photon\%inside} becomes
\texttt{.false.}), the photon weight is accumulated into the escape spectrum:
\begin{equation}
  J_\mathrm{out}[i_x] \mathrel{+}= w,\qquad
  i_x = \left\lfloor \frac{x - x_\mathrm{min}}{\delta x}\right\rfloor + 1,
\end{equation}
where $x$ is the photon frequency (global frame, in units of $\Delta\nu_D$)
and $w$ is the photon weight.  This mirrors the behaviour of
\texttt{raytrace\_to\_tau\_car}, where Jout is accumulated upon grid exit.
There are three exit paths (clump-then-sphere exit, already-outside-sphere,
and no-clump-found), all handled identically.

\subsection{Optical Depth to Edge}

\texttt{raytrace\_to\_edge\_clump} computes $\tau$ from the current position
to the sphere boundary without moving the photon (read-only input).  The same
DDA logic is used; frequency shifts at clump entry/exit are tracked locally.
This is called by peel-off routines and by the forced-first-scatter
initialisation in \texttt{run\_simulation\_mod}.

%=============================================================================
\section{Peel-Off Observations}
%=============================================================================

Peel-off is handled by the standard \texttt{peelingoff\_rect} routines
(\texttt{peeling\_direct\_outside}, \texttt{peeling\_resonance\_nostokes\_outside},
etc.) with \texttt{raytrace\_to\_edge} pointing to
\texttt{raytrace\_to\_edge\_clump}.

The observer distance defaults to $d = 100\,R$ and the image pixel size is
set so that the full sphere subtends the image:
\begin{equation}
  \delta\theta = \frac{1}{N_x/2}\arcsin\!\left(\frac{R}{d}\right)\;\mathrm{[rad]}
  \approx \frac{R}{d\,(N_x/2)}\;\mathrm{[rad]}.
\end{equation}
For $R=1$, $d=100$, $N_x=129$: $\delta\theta \approx 0\fdg0089$, covering
$\pm 1\fdg14$, which exactly spans the sphere.

%=============================================================================
\section{FITS Output of Clump Information}
%=============================================================================

When \texttt{par\%save\_clump\_info = .true.}, the subroutine
\texttt{write\_clumps\_fits} (in \texttt{clump\_mod.f90}) writes a binary-table
FITS file \texttt{<base>\_clumps.fits.gz} containing:

\paragraph{BinTable columns:}
\texttt{X}, \texttt{Y}, \texttt{Z} (clump centres, code units);
\texttt{VX}, \texttt{VY}, \texttt{VZ} (bulk velocities, km/s).

\paragraph{Header keywords (in the BinTable HDU):}
\begin{center}
\begin{tabular}{lll}
\hline
Keyword & Value & Description \\
\hline
\texttt{N\_CLUMPS} & $N$ & number of clumps (realised) \\
\texttt{SPHERE\_R} & $R$ & outer sphere radius (code units) \\
\texttt{CL\_RAD}   & $r_c$ & clump radius (code units) \\
\texttt{F\_VOL}    & $f_\mathrm{vol}$ & volume filling factor (realised) \\
\texttt{N\_COV}    & $N_\mathrm{cov}$ & covering factor (realised) \\
\texttt{TAU0}      & $\tau_0$ & line-centre $\tau$ per clump \\
\texttt{SIGMA\_V}  & $\sigma_v$ & bulk velocity dispersion (km/s) \\
\texttt{TEMP\_CL}  & $T_c$ & clump temperature used (K) \\
\texttt{RHOKAP}    & $\kappa_c$ & opacity per code unit \\
\texttt{CL\_DFREQ} & $\Delta\nu_D^{cl}$ & Doppler frequency (Hz) \\
\texttt{VTHERM}    & $v_\mathrm{th}$ & thermal velocity (km/s) \\
\texttt{VOIGT\_A}  & $a_V$ & Voigt damping parameter \\
\texttt{RMAX}      & $R$ & \texttt{par\%rmax} (input) \\
\texttt{IN\_NCOV}  & \texttt{par\%clump\_N\_cov} & input covering factor \\
\texttt{IN\_FVOL}  & \texttt{par\%clump\_f\_vol} & input filling factor \\
\texttt{IN\_NCL}   & \texttt{par\%clump\_N\_clumps} & input $N$ \\
\texttt{IN\_NH}    & \texttt{par\%clump\_nH} & input $n_H$ (cm$^{-3}$) \\
\texttt{IN\_TEMP}  & \texttt{par\%clump\_temperature} & input $T_c$ (K) \\
\hline
\end{tabular}
\end{center}

Note: column data are written first (which creates the BinTable HDU);
header keywords are then written into that same HDU.  FITS keyword names
must not exceed 8 characters.

%=============================================================================
\section{Input Parameters Summary}
%=============================================================================

\begin{center}
\begin{tabular}{lll}
\hline
Parameter & Default & Description \\
\hline
\texttt{par\%use\_clump\_medium} & \texttt{.false.} & enable clump mode \\
\texttt{par\%rmax}              & ---  & outer sphere radius (code units) \\
\texttt{par\%clump\_radius}     & ---  & clump radius $r_c$ (code units) \\
\texttt{par\%clump\_N\_cov}     & 0    & covering factor (input) \\
\texttt{par\%clump\_f\_vol}     & 0    & volume filling factor (input) \\
\texttt{par\%clump\_N\_clumps}  & 0    & number of clumps (input, direct) \\
\texttt{par\%clump\_tau0}       & 0    & line-centre $\tau$ per clump \\
\texttt{par\%clump\_nH}         & 0    & clump $n_H$ (cm$^{-3}$; requires \texttt{distance\_unit}) \\
\texttt{par\%clump\_temperature}& $-1$ & clump $T$ (K); $<0$ $\Rightarrow$ \texttt{par\%temperature} \\
\texttt{par\%clump\_sigma\_v}   & 0    & random velocity $\sigma$ per component (km/s) \\
\texttt{par\%velocity\_type}    & \texttt{''} & systematic flow type (see \S2.4) \\
\texttt{par\%Vexp}              & 0    & expansion speed for \texttt{hubble}/\texttt{constant\_radial} (km/s) \\
\texttt{par\%Vx, Vy, Vz}        & 0    & bulk velocity for \texttt{parallel\_velocity} (km/s) \\
\texttt{par\%Vpeak}             & 0    & peak speed for \texttt{ssh} (km/s) \\
\texttt{par\%rpeak}             & 0    & peak radius for \texttt{ssh} (code units) \\
\texttt{par\%DeltaV}            & 0    & outer velocity increment for \texttt{ssh} (km/s) \\
\texttt{par\%Vrot}              & 0    & rotation speed for rotating models (km/s) \\
\texttt{par\%rinner}            & 0    & solid-body radius for \texttt{rotating\_galaxy\_halo} \\
\texttt{par\%save\_clump\_info} & \texttt{.false.} & write clump FITS file \\
\hline
\end{tabular}
\end{center}

\paragraph{Example input:}
\begin{lstlisting}
&parameters
 par%use_clump_medium = .true.
 par%rmax             = 1.0
 par%clump_radius     = 0.01
 par%clump_N_cov      = 5.0
 par%clump_tau0       = 1000.0
 par%temperature      = 1.0e4
 par%clump_sigma_v    = 0.0
 par%no_photons       = 1e5
 par%spectral_type    = 'voigt'
 par%nxfreq           = 301
 par%velocity_min     = -100.0
 par%velocity_max     =  100.0
 par%nx = 11, par%ny = 11, par%nz = 11
 par%nxim = 129, par%nyim = 129
 par%save_clump_info  = .true.
 par%out_file         = 'clump_test.fits.gz'
/
\end{lstlisting}

%=============================================================================
\section{Memory Layout}
%=============================================================================

Clump positions, velocities, and CSR arrays are allocated as MPI-3
shared-memory windows via \texttt{create\_shared\_mem} from
\texttt{memory\_mod}, giving one copy per compute node rather than per MPI
rank.  On systems with multiple ranks per node this significantly reduces
memory usage for large clump populations.

Cleanup is performed by \texttt{destroy\_clumps} which calls
\texttt{destroy\_shared\_mem\_all} (collective \texttt{MPI\_WIN\_FREE} across
all ranks) followed by explicit \texttt{nullify} of all pointers.  The
individual \texttt{destroy\_mem} is not used because
\texttt{MPI\_WIN\_SHARED\_QUERY} returns the rank-0 base address, which
differs from the local pointer on non-root ranks.

%=============================================================================
\section{Validation and Examples}
%=============================================================================

The directory \texttt{examples/clump\_sphere/} contains a suite of six
reference runs (plus the legacy \texttt{clump\_test.in}) covering the
main parameter space.  All named cases use $R=1$, $r_c=0.01$, $10^6$
photons, and \texttt{par\%use\_clump\_medium = .true.}:

\begin{center}
\begin{tabular}{lll}
\hline
Input file & $\tau_0$ & Notes \\
\hline
\texttt{clump\_tau3\_ncov5.in}        & $10^3$ & $N_\mathrm{cov}=5$, static (reference) \\
\texttt{clump\_tau4\_ncov5.in}        & $10^4$ & $N_\mathrm{cov}=5$, static \\
\texttt{clump\_tau3\_ncov20.in}       & $10^3$ & $N_\mathrm{cov}=20$, dense covering \\
\texttt{clump\_tau3\_ncov5\_hubble.in}& $10^3$ & $N_\mathrm{cov}=5$, Hubble $V_\mathrm{exp}=200$ km/s \\
\texttt{clump\_tau3\_ncov5\_sigv.in}  & $10^3$ & $N_\mathrm{cov}=5$, $\sigma_v=50$ km/s \\
\texttt{clump\_tau3\_ncov5\_peel.in}  & $10^3$ & $N_\mathrm{cov}=5$, peel-off $129\times129$ \\
\hline
\end{tabular}
\end{center}

Run with \texttt{bash run.sh} (uses 4 MPI ranks by default) or individually:
\begin{verbatim}
mpirun -np 4 ../../LaRT.x clump_tau3_ncov5.in
\end{verbatim}
Python post-processing:
\begin{verbatim}
python plot_clump.py          # all runs
python plot_clump.py clump_tau3_ncov5   # single run
\end{verbatim}

For the reference case ($N_\mathrm{cov}=5$, $\tau_0=10^3$):
\begin{itemize}
  \item Realised $N_\mathrm{cov} \approx 5.00$, $f_\mathrm{vol} \approx 0.067$;
  \item RSA acceptance rate $\sim 74\%$, $N=66667$ clumps placed;
  \item Mean scattering count $\approx 6400$ per photon;
  \item Wall time $\approx 3.75$ min (10 ranks, $10^5$ photons).
\end{itemize}

The escaped spectrum $J_\mathrm{out}(x)$ shows the characteristic double-peak
expected for a clumpy Lyman-$\alpha$ source.  The peel-off output
(\texttt{\_obs.fits.gz}) contains \texttt{Scattered} (HDU~0) and
\texttt{Direct} (HDU~1) spectral image cubes of shape $(301, 129, 129)$.

\end{document}
